% rainbow-test.tex
% rainbow.sty 测试文档 —— 展示所有定理类环境
\documentclass[12pt,a4paper]{ctexart}
\usepackage{amsmath,amssymb}
\usepackage{rainbow}
\usepackage{geometry}
\geometry{margin=2.5cm}

\title{\textbf{rainbow.sty 环境测试}}
\author{测试文档}
\date{\today}

\begin{document}
\maketitle
\tableofcontents
\newpage

%% ============================================================
\section{基本定理环境}
%% ============================================================

\begin{definition}[拓扑空间]
设 $X$ 是一个非空集合，$\mathcal{T}\subseteq\mathcal{P}(X)$ 是 $X$ 的子集族。
若 $\mathcal{T}$ 满足以下三个条件：
\begin{enumerate}
  \item $\varnothing,\,X\in\mathcal{T}$；
  \item $\mathcal{T}$ 中任意多个成员的并仍属于 $\mathcal{T}$；
  \item $\mathcal{T}$ 中有限多个成员的交仍属于 $\mathcal{T}$，
\end{enumerate}
则称 $\mathcal{T}$ 是 $X$ 上的一个\textbf{拓扑}，$(X,\mathcal{T})$ 称为\textbf{拓扑空间}。
\end{definition}

\begin{definition}[连续映射]
设 $(X,\mathcal{T}_X)$ 和 $(Y,\mathcal{T}_Y)$ 是拓扑空间，映射 $f\colon X\to Y$。
若对任意 $V\in\mathcal{T}_Y$，都有 $f^{-1}(V)\in\mathcal{T}_X$，
则称 $f$ 是\textbf{连续}的。
\end{definition}

\begin{theorem}[Bolzano--Weierstrass]
$\mathbb{R}^n$ 中的有界无穷子集必有聚点。
\end{theorem}

\begin{proof}
利用闭区间套定理和逐次二分法即可证明。
\end{proof}

\begin{lemma}[Urysohn 引理]
设 $X$ 是正规空间，$A,B$ 是 $X$ 中两个不相交的闭集，
则存在连续函数 $f\colon X\to[0,1]$ 使得
$f|_A=0$，$f|_B=1$。
\end{lemma}

\begin{corollary}
正规空间中，任意两个不相交闭集可以被连续函数分离。
\end{corollary}

\begin{theorem}[Heine--Borel]
$\mathbb{R}^n$ 中的子集是紧致的，当且仅当它是有界闭集。
\end{theorem}

\begin{corollary}[有限覆盖]
$[a,b]$ 上的每个开覆盖都有有限子覆盖。
\end{corollary}

%% ============================================================
\section{更多环境展示}
%% ============================================================

\begin{axiom}[Zermelo 选择公理]
设 $\{A_i\}_{i\in I}$ 是一族非空集合，则存在一个选择函数
$f\colon I\to\bigcup_{i\in I}A_i$，使得对每个 $i\in I$，$f(i)\in A_i$。
\end{axiom}

\begin{postulate}[欧几里得第五公设]
过直线外一点，有且仅有一条直线与已知直线平行。
\end{postulate}

\begin{proposition}[开映射性质]
设 $f\colon X\to Y$ 是连续双射且 $X$ 紧致、$Y$ Hausdorff，
则 $f$ 是同胚。
\end{proposition}

\begin{proof}
只需证 $f$ 是闭映射。紧致空间中的闭集是紧致的，紧致集在连续映射下的像是紧致的，
而 Hausdorff 空间中的紧致子集是闭的，因此 $f$ 将闭集映为闭集。
\end{proof}

\begin{assumption}[有界性假设]
假设对所有 $n\geq1$，存在常数 $C>0$ 使得 $\|x_n\|\leq C$。
\end{assumption}

\begin{property}[Banach 空间的完备性]
Banach 空间中的每个 Cauchy 列都收敛。
\end{property}

\begin{conjecture}[Goldbach 猜想]
每个大于 $2$ 的偶数都可以表示为两个素数之和。
\end{conjecture}

%% ============================================================
\section{编号重置验证}
%% ============================================================

进入新的 section 后，所有计数器应当从 $1$ 重新开始。

\begin{theorem}
本节第一个定理，编号应为 3.1。
\end{theorem}

\begin{definition}
本节第一个定义，编号应为 3.1。
\end{definition}

\begin{lemma}
本节第一个引理，编号应为 3.1。
\end{lemma}

\begin{proposition}
本节第一个命题，编号应为 3.1。
\end{proposition}

\begin{corollary}
本节第一个推论，编号应为 3.1。
\end{corollary}

\begin{theorem}[第二个定理]
编号应为 3.2，验证定理计数器独立递增。
\end{theorem}

%% ============================================================
\section{remark 和 example（无编号）}
%% ============================================================

\begin{remark}
注记环境没有编号，使用灰色左边框。
\end{remark}

\begin{example}
这是一个例子环境，同样没有编号。
考虑 $f(x)=x^2$ 在 $[0,1]$ 上的积分：
\[
  \int_0^1 x^2\,dx = \frac{1}{3}.
\]
\end{example}

\end{document}
